% Источники: exam/materials/преподаватель/2026/Моделирование_случайных_величин,_векторов_.pdf; exam/materials/моделирование.pdf, раздел 13.
\section{Моделирование случайных величин и векторов с заданными законами распределения.}

\subsection*{Источник случайности}

Аппаратный или программный генератор псевдослучайных чисел вырабатывает последовательность $r_1, r_2, \ldots$ из равномерного распределения $U(0,1)$.

\subsection*{Метод обращения (инверсии)}

Если $W(x)$ — функция распределения (ФР) моделируемой СВ $\xi$, то
\[
  \xi = W^{-1}(r), \quad r \sim U(0,1).
\]
Применим при наличии явной аналитической формулы $W^{-1}$.

Примеры: экспоненциальный закон $W(x) = 1 - e^{-\lambda x}$, $\xi = -\frac{1}{\lambda}\ln(1-r)$.

\subsection*{Метод Неймана (отбора-отклонения)}

Применяется, когда $W^{-1}$ не имеет удобного аналитического выражения.
\begin{enumerate}
  \item Выбирают вспомогательное распределение $g(x)$ и константу $M \geqslant f(x)/g(x)$ для всех $x$.
  \item Генерируют $x \sim g(x)$ и $r \sim U(0,1)$.
  \item Принять $\xi = x$, если $r \leqslant f(x)/(M\,g(x))$; иначе повторить с шага 1.
\end{enumerate}

\subsection*{Моделирование случайных векторов}

Для вектора $\boldsymbol{\xi} = (\xi_1, \ldots, \xi_m)^\top$ с совместной функцией распределения $F(\mathbf{x})$ используют \textbf{разложение на условные распределения}:
\[
  f(x_1,\ldots,x_m) = f_1(x_1)\, f_{2|1}(x_2|x_1)\cdots f_{m|1\ldots m-1}(x_m|x_1,\ldots,x_{m-1}).
\]
Компоненты генерируются последовательно: сначала $\xi_1$ по $F_1$, затем $\xi_2$ по условному распределению $F_{2|1}(\cdot|\xi_1)$ и так далее до $\xi_m$ по $F_{m|1\ldots m-1}(\cdot|\xi_1,\ldots,\xi_{m-1})$.

Для нормального вектора применяется \textbf{разложение Холецкого}: $\boldsymbol{\xi} = \boldsymbol{\mu} + A\mathbf{z}$, где $A$ — нижнетреугольная матрица из $\Sigma = AA^\top$, $\mathbf{z}$ — вектор независимых $N(0,1)$.

\clearpage
